\part{Eight-Vertex / RSOS Models}

\section{Theta and elliptic functions}

\begin{exercise}[Restricted solid-on-solid configurations]
Fix an integer $k\ge1$. The height set of the restricted solid-on-solid
(RSOS) model is $\{1,\dots,k+1\}$; neighbouring lattice sites must
have heights differing by one. On a finite graph take the uniform
measure on admissible configurations, with zero weight for
inadmissible ones. Equivalently the energy is zero on the allowed
configuration space.
Prove invariance under the height reflection $a\mapsto k+2-a$.
On a path of $N$ edges with initial height $a$ and terminal height
$b$, express the partition function as $(A^N)_{ab}$, where $A$
is the adjacency kernel of the path graph $A_{k+1}$.
\end{exercise}

\begin{exercise}[Jacobi theta functions: series, quasi-periodicity, half-period shifts, and zeros]
Let $\tau\in\C$ with $\operatorname{Im}\tau>0$ and $p=e^{\ii\pi\tau}$ (the
nome), so $|p|<1$; fix $p^{1/4}=e^{\ii\pi\tau/4}$.  Define, for $z\in\C$,
\begin{align*}
 \vartheta_1(z|\tau)&=2\sum_{n\ge0}(-1)^np^{(n+1/2)^2}\sin((2n+1)z)
 =-\ii\sum_{n\in\Z}(-1)^np^{(n+1/2)^2}e^{\ii(2n+1)z},\\
 \vartheta_2(z|\tau)&=2\sum_{n\ge0}p^{(n+1/2)^2}\cos((2n+1)z),\qquad
 \vartheta_3(z|\tau)=1+2\sum_{n\ge1}p^{n^2}\cos(2nz),\\
 \vartheta_4(z|\tau)&=1+2\sum_{n\ge1}(-1)^np^{n^2}\cos(2nz)=\sum_{n\in\Z}(-1)^np^{n^2}e^{2\ii nz}.
\end{align*}
The dependence on $\tau$ is suppressed when it is fixed.
\begin{enumerate}[label=(\alph*)]
\item Prove that the series converge absolutely and uniformly on compact
sets, so the $\vartheta_j$ are entire; $\vartheta_1$ is odd and
$\vartheta_2,\vartheta_3,\vartheta_4$ are even.
\item Prove the quasi-periodicity relations
\begin{align*}
 \vartheta_1(z+\pi)&=-\vartheta_1(z),&\vartheta_1(z+\pi\tau)&=-p^{-1}e^{-2\ii z}\vartheta_1(z),\\
 \vartheta_2(z+\pi)&=-\vartheta_2(z),&\vartheta_2(z+\pi\tau)&=p^{-1}e^{-2\ii z}\vartheta_2(z),\\
 \vartheta_3(z+\pi)&=\vartheta_3(z),&\vartheta_3(z+\pi\tau)&=p^{-1}e^{-2\ii z}\vartheta_3(z),\\
 \vartheta_4(z+\pi)&=\vartheta_4(z),&\vartheta_4(z+\pi\tau)&=-p^{-1}e^{-2\ii z}\vartheta_4(z),
\end{align*}
and the half-period shifts
$\vartheta_1(z+\tfrac\pi2)=\vartheta_2(z)$, $\vartheta_2(z+\tfrac\pi2)=-\vartheta_1(z)$,
$\vartheta_3(z+\tfrac\pi2)=\vartheta_4(z)$, $\vartheta_4(z+\tfrac\pi2)=\vartheta_3(z)$,
$\vartheta_1(z+\tfrac{\pi\tau}2)=\ii p^{-1/4}e^{-\ii z}\vartheta_4(z)$,
$\vartheta_4(z+\tfrac{\pi\tau}2)=\ii p^{-1/4}e^{-\ii z}\vartheta_1(z)$.
\item Zero count.  Let $f$ be entire, not identically zero, with
$f(z+\pi)=\epsilon f(z)$, $\epsilon\in\{\pm1\}$, and
$f(z+\pi\tau)=Ce^{-2\ii nz}f(z)$ for some $C\in\C^\times$ and $n\in\Z_{\ge0}$.
Prove by integrating $f'/f$ around a period parallelogram avoiding zeros
that $f$ has exactly $n$ zeros modulo the lattice $\pi\Z+\pi\tau\Z$,
counted with multiplicity.  Deduce that $\vartheta_1$ vanishes exactly on
the lattice, with simple zeros, and that the zeros of
$\vartheta_2,\vartheta_3,\vartheta_4$ are the translates of the lattice by
$\frac\pi2$, $\frac\pi2+\frac{\pi\tau}2$, $\frac{\pi\tau}2$ respectively.
\item Dimension count.  For $\epsilon\in\{\pm1\}$, $n\ge1$, $C\in\C^\times$
let $\Theta_n(\epsilon,C)$ be the space of entire $f$ with
$f(z+\pi)=\epsilon f(z)$ and $f(z+\pi\tau)=Ce^{-2\ii nz}f(z)$.  Expanding
$f(z)=\sum_mc_me^{\ii(2m+\delta)z}$ with $\delta\in\{0,1\}$ determined by
$\epsilon$, show that the second relation is a recursion $c_{m+n}=\gamma_mc_m$
with explicit $\gamma_m\ne0$, whose solutions converge, and conclude
$\dim\Theta_n(\epsilon,C)=n$.  Verify that a product of $n$ theta
functions $\vartheta_{j_1}(z+w_1)\cdots\vartheta_{j_n}(z+w_n)$ lies in
$\Theta_n(\epsilon,C)$ for suitable $\epsilon,C$ depending only on the
$j_i$ and $\sum w_i$.
\end{enumerate}
\end{exercise}

\begin{exercise}[Addition identities of degree two]
Write $\vartheta_j=\vartheta_j(\cdot|\tau)$.  Prove each of the following
identities by the same method: both sides lie in a space
$\Theta_2(\epsilon,C)$ of dimension two (as functions of $x$, with the
other variables fixed), the right side vanishes at a point where the left
side also vanishes, so that the two sides are proportional, and the constant
is found by evaluating at $x=0$ or $x=\frac\pi2$.
\begin{enumerate}[label=(\alph*)]
\item $\vartheta_1(x+y)\vartheta_1(x-y)\vartheta_4(0)^2=\vartheta_1(x)^2\vartheta_4(y)^2-\vartheta_4(x)^2\vartheta_1(y)^2$
and $\vartheta_4(x+y)\vartheta_4(x-y)\vartheta_4(0)^2=\vartheta_4(x)^2\vartheta_4(y)^2-\vartheta_1(x)^2\vartheta_1(y)^2$.
\item $\vartheta_1(x+y)\vartheta_4(x-y)\vartheta_2(0)\vartheta_3(0)=\vartheta_1(x)\vartheta_4(x)\vartheta_2(y)\vartheta_3(y)+\vartheta_2(x)\vartheta_3(x)\vartheta_1(y)\vartheta_4(y)$.
\item $\vartheta_4(0)^2\vartheta_3(x)^2=\vartheta_3(0)^2\vartheta_4(x)^2-\vartheta_2(0)^2\vartheta_1(x)^2$
and $\vartheta_4(0)^2\vartheta_2(x)^2=\vartheta_2(0)^2\vartheta_4(x)^2-\vartheta_3(0)^2\vartheta_1(x)^2$;
deduce $\vartheta_3(0)^4=\vartheta_2(0)^4+\vartheta_4(0)^4$ by putting $x=\frac\pi2$.
\item For a fixed shift $\eta$ put
$f(x)=\vartheta_4(x)\vartheta_1(x+\eta)$,
$g(x)=\vartheta_1(x)\vartheta_4(x+\eta)$,
$h(x)=\vartheta_4(x)\vartheta_4(x+\eta)$, and
$j(x)=\vartheta_1(x)\vartheta_1(x+\eta)$.
Prove the four identities
\begin{align*}
 f(x)f(y)-g(x)g(y)
 &=\vartheta_4(0)\vartheta_1(\eta)
       \vartheta_4(x-y)\vartheta_1(x+y+\eta),\\
 h(x)h(y)-j(x)j(y)
 &=\vartheta_4(0)\vartheta_4(\eta)
       \vartheta_4(x-y)\vartheta_4(x+y+\eta),\\
 h(x)j(y)-j(x)h(y)
 &=\vartheta_4(0)\vartheta_4(\eta)
       \vartheta_1(y-x)\vartheta_1(x+y+\eta),\\
 f(x)g(y)-g(x)f(y)
 &=\vartheta_4(0)\vartheta_1(\eta)
       \vartheta_1(y-x)\vartheta_4(x+y+\eta).
\end{align*}
For the first use the common zero $x=-y-\eta$; for the
second use $x=y+\pi\tau/2$ and the half-period shifts;
for the last two use $x=y$.
\item The three-term identity.  Put $[x,y]=\vartheta_1(x+y)\vartheta_1(x-y)$.
Deduce from (a) that $[x,y]\vartheta_4(0)^2=P_xQ_y-Q_xP_y$ with
$P_x=\vartheta_1(x)^2$, $Q_x=\vartheta_4(x)^2$, and from the Pl\"ucker
relation for $2\times2$ determinants that
\[
 [c,d]\,[a,b]+[d,b]\,[a,c]+[b,c]\,[a,d]=0\qquad(a,b,c,d\in\C).
\]
\item Wronskian.  Show that $\vartheta_1'\vartheta_4-\vartheta_1\vartheta_4'$
lies in the same $\Theta_2$ as $\vartheta_2\vartheta_3$ and is even, and
conclude $(\vartheta_1/\vartheta_4)'=\kappa\,\vartheta_2\vartheta_3/\vartheta_4^2$
with $\kappa=\vartheta_1'(0)\vartheta_4(0)/(\vartheta_2(0)\vartheta_3(0))$.
\end{enumerate}
\end{exercise}

\begin{exercise}[Jacobi's triple product and the product formulas]
\leavevmode\par
\begin{enumerate}[label=(\alph*)]
\item For $|p|<1$ and $w\in\C^\times$ put
$F(w)=\prod_{n\ge1}(1+p^{2n-1}w)(1+p^{2n-1}w^{-1})$.  Prove that $F$ is
holomorphic on $\C^\times$ with $F(p^2w)=(pw)^{-1}F(w)$, expand
$F(w)=\sum_{m\in\Z}c_mw^m$, derive $c_m=p^{2m-1}c_{m-1}$, hence
$c_m=p^{m^2}c_0$, so that $\prod_{n\ge1}(1-p^{2n})\,F(w)=\phi(p)\sum_mp^{m^2}w^m$
for some $\phi(p)$ depending only on $p$.
\item Determine $\phi(p)=1$.  Put $w=\ii$: the series becomes
$\sum_j(-1)^jp^{4j^2}$ and $F(\ii)=\prod_n(1+p^{4n-2})$.  Apply (a) with
$p$ replaced by $p^4$ and $w=-1$ to the same series, and use
$\prod_n(1-p^{2n})=\prod_n(1-p^{8n})(1-p^{8n-4})(1-p^{4n-2})$ to obtain
$\phi(p)=\phi(p^4)$; iterate and let $p\to0$.
\item Deduce the product formulas
\begin{align*}
 \vartheta_3(z)&=\prod_{n\ge1}(1-p^{2n})(1+2p^{2n-1}\cos2z+p^{4n-2}),\\
 \vartheta_4(z)&=\prod_{n\ge1}(1-p^{2n})(1-2p^{2n-1}\cos2z+p^{4n-2}),\\
 \vartheta_2(z)&=2p^{1/4}\cos z\prod_{n\ge1}(1-p^{2n})(1+2p^{2n}\cos2z+p^{4n}),\\
 \vartheta_1(z)&=2p^{1/4}\sin z\prod_{n\ge1}(1-p^{2n})(1-2p^{2n}\cos2z+p^{4n}),
\end{align*}
and the derivative formula $\vartheta_1'(0)=\vartheta_2(0)\vartheta_3(0)\vartheta_4(0)$
from $\prod_n(1+p^{2n})(1+p^{2n-1})(1-p^{2n-1})=1$.  Conclude that the
constant $\kappa$ of the previous exercise equals $\vartheta_4(0)^2$.
\item Degeneration.  Show that as $p\to0$ with $z$ fixed,
$\vartheta_1(z)=2p^{1/4}\sin z\,(1+O(p^2))$, $\vartheta_2(z)=2p^{1/4}\cos z\,(1+O(p^2))$,
$\vartheta_3(z),\vartheta_4(z)=1+O(p)$, uniformly on compact sets.
\end{enumerate}
\end{exercise}

\begin{exercise}[Modular transformation, Jacobi elliptic functions, and the modulus]
\leavevmode\par
\begin{enumerate}[label=(\alph*)]
\item Prove the Fourier identity
$\sum_{n\in\Z}e^{-\pi n^2t+2\pi\ii nz}=t^{-1/2}\sum_{n\in\Z}e^{-\pi(n-z)^2/t}$
for $\operatorname{Re}t>0$, $z\in\C$ by integrating a Gaussian against $e^{-2\pi\ii nx}$ and
periodizing the result (the branch of $t^{-1/2}$ is positive
for $t>0$). Derive
\[
 \vartheta_3(z|\tau)=(-\ii\tau)^{-1/2}e^{-\ii z^2/(\pi\tau)}\,\vartheta_3\!\Big(\frac z\tau\Big|-\frac1\tau\Big),
\]
and, using the half-period shifts,
$\vartheta_4(z|\tau)=(-\ii\tau)^{-1/2}e^{-\ii z^2/(\pi\tau)}\vartheta_2(z/\tau|-1/\tau)$,
$\vartheta_2(z|\tau)=(-\ii\tau)^{-1/2}e^{-\ii z^2/(\pi\tau)}\vartheta_4(z/\tau|-1/\tau)$,
$\vartheta_1(z|\tau)=\ii(-\ii\tau)^{-1/2}e^{-\ii z^2/(\pi\tau)}\vartheta_1(z/\tau|-1/\tau)$.
\item Define the modulus $k=\vartheta_2(0)^2/\vartheta_3(0)^2$, the
complementary modulus $k'=\vartheta_4(0)^2/\vartheta_3(0)^2$, and
$K=\frac\pi2\vartheta_3(0)^2$, $K'=-\ii\tau K$.  Show $k^2+k'^2=1$ and that
$\tau\mapsto-1/\tau$ exchanges $k$ with $k'$ and $K$ with $K'$.
\item Define, with $u=\vartheta_3(0)^2z$,
\[
 \operatorname{sn}u=\frac{\vartheta_3(0)}{\vartheta_2(0)}\frac{\vartheta_1(z)}{\vartheta_4(z)},\qquad
 \operatorname{cn}u=\frac{\vartheta_4(0)}{\vartheta_2(0)}\frac{\vartheta_2(z)}{\vartheta_4(z)},\qquad
 \operatorname{dn}u=\frac{\vartheta_4(0)}{\vartheta_3(0)}\frac{\vartheta_3(z)}{\vartheta_4(z)}.
\]
Prove from the identities of degree two that
$\operatorname{sn}^2+\operatorname{cn}^2=1$, $\operatorname{dn}^2+k^2\operatorname{sn}^2=1$,
$\frac{\dd}{\dd u}\operatorname{sn}=\operatorname{cn}\operatorname{dn}$, that
$\operatorname{sn}$ is meromorphic with periods $4K$ and $2\ii K'$, simple
poles at $\ii K'+2K\Z+2\ii K'\Z$, and that $\operatorname{sn}$ is
determined by the initial value problem
$(\operatorname{sn}')^2=(1-\operatorname{sn}^2)(1-k^2\operatorname{sn}^2)$, $\operatorname{sn}0=0$, $\operatorname{sn}'0=1$,
so that it inverts the elliptic integral
$u=\int_0^{x}\dd t/\sqrt{(1-t^2)(1-k^2t^2)}$ near $0$.  For $0<k<1$
(i.e. $\tau\in\ii\R_{>0}$) show that $\operatorname{sn}$ increases from
$0$ to $1$ on $[0,K]$ (use $\operatorname{sn}K=1$ from the half-period
shift) and deduce $K=\int_0^{\pi/2}(1-k^2\sin^2\theta)^{-1/2}\dd\theta$
by the substitution $\operatorname{sn}u=\sin\theta$.
\item Show that as $p\to0$: $k\to0$, $K\to\pi/2$, $K'\to\infty$,
$\operatorname{sn}u\to\sin u$, $\operatorname{cn}u\to\cos u$, $\operatorname{dn}u\to1$,
uniformly on compact sets; and that $\tau\mapsto-1/\tau$ followed by
$p\to0$ gives the hyperbolic degeneration $\operatorname{sn}(u;k)\to\tanh u$ as $k\to1$.
\end{enumerate}
\end{exercise}

\section{Elliptic R-matrices}

\begin{exercise}[Eight-vertex configurations, weights, invariants, and the elliptic parametrization]
On the square lattice place an arrow on every edge and keep the vertices at
which the number of incoming arrows is even: the six vertices of the
section Row-to-row transfer matrices in Part IV, with weights $a$
(both arrows continue straight, parallel), $b$ (straight, antiparallel),
$c$ (turning), together with the two vertices with four incoming or four
outgoing arrows (source and sink), both of weight $d$.  With $V=\C^2$,
$e_+$ (arrow up or right), $e_-$, the vertex operator is
\[
 R=\begin{pmatrix}a&0&0&d\\0&b&c&0\\0&c&b&0\\d&0&0&a\end{pmatrix}\in\End(V\otimes V)
\]
in the basis $e_+\otimes e_+,e_+\otimes e_-,e_-\otimes e_+,e_-\otimes e_-$,
and the partition function on the $L\times N$ torus is $Z=\Tr(T^N)$ with
$T$ the row-to-row transfer matrix built from $R$ exactly as in Part IV.
Put
\[
 \Delta=\frac{a^2+b^2-c^2-d^2}{2(ab+cd)},\qquad\Gamma=\frac{ab-cd}{ab+cd}.
\]
\begin{enumerate}[label=(\alph*)]
\item Prove that $R$ commutes with $\sigma^z\otimes\sigma^z$ and with
$\sigma^x\otimes\sigma^x$ (parity of the number of down arrows is conserved,
and arrow reversal is a symmetry), and that $R$ commutes with
$\sigma^z\otimes1+1\otimes\sigma^z$ if and only if $d=0$.
\item Elliptic parametrization.  Fix $\tau$, $\eta\in\C$ and define, with
$\vartheta_j=\vartheta_j(\cdot|\tau)$,
\begin{align*}
 a(u)&=\vartheta_4(\eta)\,\vartheta_4(u)\,\vartheta_1(u+\eta),&
 b(u)&=\vartheta_4(\eta)\,\vartheta_1(u)\,\vartheta_4(u+\eta),\\
 c(u)&=\vartheta_1(\eta)\,\vartheta_4(u)\,\vartheta_4(u+\eta),&
 d(u)&=\vartheta_1(\eta)\,\vartheta_1(u)\,\vartheta_1(u+\eta),
\end{align*}
and $R(u)$ accordingly.  Show that these are Baxter's weights
\cite{baxter-eightvertex} $\vartheta_4(2\eta')\vartheta_4(u'-\eta')\vartheta_1(u'+\eta')$,
etc., under $u'=u+\eta/2$, $\eta'=\eta/2$, up to the common
factor.  Prove $\Gamma=\dfrac{\vartheta_4(\eta)^2-\vartheta_1(\eta)^2}{\vartheta_4(\eta)^2+\vartheta_1(\eta)^2}$
directly, and
\[
 \Delta=\frac{\vartheta_2(\eta)\vartheta_3(\eta)\vartheta_4(0)^2}{\vartheta_2(0)\vartheta_3(0)\,(\vartheta_4(\eta)^2+\vartheta_1(\eta)^2)}
\]
by combining the identities (a) and (b) of the exercise on addition
identities: write $a^2+b^2-c^2-d^2$ as
$\vartheta_1(u+\eta)^2[\vartheta_4(\eta)^2\vartheta_4(u)^2-\vartheta_1(\eta)^2\vartheta_1(u)^2]
+\vartheta_4(u+\eta)^2[\vartheta_4(\eta)^2\vartheta_1(u)^2-\vartheta_1(\eta)^2\vartheta_4(u)^2]$,
factor each bracket, and use (b) at $y=\pm\eta$.  Both invariants are
independent of $u$.
\item Trigonometric limit.  Show that as $p\to0$,
$(a,b,c,d)/(2p^{1/4})\to(\sin(u+\eta),\sin u,\sin\eta,0)$, which is the
six-vertex parametrization of Part IV with $(u,\eta)$ replaced by
$(\ii u,\ii\eta)$ up to a common factor, and $\Gamma\to1$,
$\Delta\to\cos\eta$.
\item Positivity.  For $\tau\in\ii\R_{>0}$ (so $0<p<1$) show from the
product formulas that $\vartheta_4(x)>0$ for real $x$ and
$\vartheta_1(x)>0$ for $0<x<\pi$; deduce that for real $\eta,u$ with
$0<u$ and $u+\eta<\pi$, $0<\eta<\pi$, all four weights are positive, and
compute the signs of $\Delta$ and $\Gamma$ there.
\end{enumerate}
\end{exercise}

\begin{exercise}[The eight-vertex Yang--Baxter equation]
Let $R(u)$ be the elliptic eight-vertex $R$-matrix.
\begin{enumerate}[label=(\alph*)]
\item Write the equation
$R_{12}(u-v)R_{13}(u)R_{23}(v)=R_{23}(v)R_{13}(u)R_{12}(u-v)$ on
$V^{\otimes3}$ with three sets of weights $(a,b,c,d)$, $(a',b',c',d')$,
$(a'',b'',c'',d'')$ for the arguments $u-v$, $u$, $v$.  Using the
parity conservation and arrow reversal of the previous exercise, show that
it is equivalent to the six polynomial identities
\begin{align*}
 d(a'a''-b'b'')&=a(c''d'-c'd''),&
 c(a'b''-a''b')&=b(d'd''-c'c''),\\
 c'(aa''-bb'')&=a'(cc''-dd''),&
 c''(ab'-a'b)&=b''(cc'-dd'),\\
 d''(aa'-bb')&=a''(cd'-c'd),&
 d'(ab''-a''b)&=b'(cd''-c''d).
\end{align*}
\item Using the four shifted identities of the exercise on addition
identities, prove for all $x,y$:
\begin{align*}
 a(x)a(y)-b(x)b(y)&=\vartheta_4(\eta)^2\vartheta_4(0)\vartheta_1(\eta)\,\vartheta_4(x-y)\vartheta_1(x+y+\eta),\\
 c(x)c(y)-d(x)d(y)&=\vartheta_1(\eta)^2\vartheta_4(0)\vartheta_4(\eta)\,\vartheta_4(x-y)\vartheta_4(x+y+\eta),\\
 c(x)d(y)-c(y)d(x)&=\vartheta_1(\eta)^2\vartheta_4(0)\vartheta_4(\eta)\,\vartheta_1(y-x)\vartheta_1(x+y+\eta),\\
 a(x)b(y)-a(y)b(x)&=\vartheta_4(\eta)^2\vartheta_4(0)\vartheta_1(\eta)\,\vartheta_1(y-x)\vartheta_4(x+y+\eta).
\end{align*}
\item Substitute $(x,y)=(u-v,v)$, $(u-v,u)$, $(u,v)$ as appropriate into
(b) and verify each of the six identities of (a); for instance the third
becomes
$c(u)\,\vartheta_4(\eta)^2\vartheta_4(0)\vartheta_1(\eta)\vartheta_4(u-2v)\vartheta_1(u+\eta)
=a(u)\,\vartheta_1(\eta)^2\vartheta_4(0)\vartheta_4(\eta)\vartheta_4(u-2v)\vartheta_4(u+\eta)$,
which holds by the definitions of $a(u)$ and $c(u)$.  Conclude that the
eight-vertex Yang--Baxter equation holds identically in $u,v,\eta,\tau$.
\item Deduce the three-parameter form
$R_{12}(u-v)R_{13}(u-w)R_{23}(v-w)=R_{23}(v-w)R_{13}(u-w)R_{12}(u-v)$
and, in the limit $p\to0$, recover the six-vertex Yang--Baxter equation
of Part IV.
\end{enumerate}
\end{exercise}

\begin{exercise}[Regularity, unitarity, crossing symmetry, and commuting eight-vertex transfer matrices]
\leavevmode\par
\begin{enumerate}[label=(\alph*)]
\item Prove $R(0)=a(0)P$ with $a(0)=\vartheta_4(\eta)\vartheta_4(0)\vartheta_1(\eta)$
and $P$ the flip (regularity).
\item Prove $R_{12}(u)R_{12}(-u)=\rho(u)\,1$ with
$\rho(u)=a(u)a(-u)+d(u)d(-u)=b(u)b(-u)+c(u)c(-u)$ (unitarity; use the
identities of the previous exercise at $(x,y)=(u,-u)$ to show the two
expressions agree and the cross terms
\[
 a(u)d(-u)+d(u)a(-u),\qquad b(u)c(-u)+c(u)b(-u)
\]
vanish).
\item Prove the crossing relations $a(-u-\eta)=-b(u)$, $b(-u-\eta)=-a(u)$,
$c(-u-\eta)=c(u)$, $d(-u-\eta)=d(u)$, and deduce
\[
 R_{12}(-u-\eta)=-(\sigma^y\otimes1)R_{12}(u)^{t_1}
                            (\sigma^y\otimes1)
\]
where $t_1$ is the partial transpose in the first factor (crossing symmetry).
\item Define $L_j(u)=R_{0j}(u)$, $\mathcal T(u)=L_N(u)\cdots L_1(u)$,
$t(u)=\Tr_0\mathcal T(u)$ on $\mathcal H_N$.  Prove the RLL relation and
$[t(u),t(v)]=0$ exactly as in the section Commuting families in Part IV,
now without conservation of $S^z$ but with $[t(u),\sigma^z_1\cdots\sigma^z_N]=0$
and $[t(u),\sigma^x_1\cdots\sigma^x_N]=0$; show that the row-to-row
transfer matrix of the eight-vertex model on the torus is $t(u)$ for the
appropriate normalization, so that on an $N$-site row and $M$ rows, $Z=\Tr t(u)^M$.
\end{enumerate}
\end{exercise}

\begin{exercise}[The XYZ chain from the logarithmic derivative and the hierarchy of degenerations]
\leavevmode\par
\begin{enumerate}[label=(\alph*)]
\item Repeating the derivation of the section Rational/trigonometric
hierarchy in Part V with $a(0)$ in place of $\sinh\eta$, prove
\[
\begin{aligned}
 t(0)^{-1}t'(0)&=\frac1{a(0)}\sum_{j=1}^NP_{j,j+1}R'_{j,j+1}(0),\\
 P R'(0)&=\frac{a'+c'}2+\frac{a'-c'}2\sigma^z\otimes\sigma^z\\
 &\quad+\frac{b'+d'}2\sigma^x\otimes\sigma^x
       +\frac{b'-d'}2\sigma^y\otimes\sigma^y .
\end{aligned}
\]
with $a'=a'(0)$ etc.
\item Compute $a'(0)=\vartheta_4(\eta)\vartheta_4(0)\vartheta_1'(\eta)$,
$b'(0)=\vartheta_4(\eta)^2\vartheta_1'(0)$, $c'(0)=\vartheta_1(\eta)\vartheta_4(0)\vartheta_4'(\eta)$,
$d'(0)=\vartheta_1(\eta)^2\vartheta_1'(0)$, and, using the Wronskian identity
$(\vartheta_1/\vartheta_4)'=\vartheta_4(0)^2\vartheta_2\vartheta_3/\vartheta_4^2$
and $\vartheta_1'(0)=\vartheta_2(0)\vartheta_3(0)\vartheta_4(0)$, prove that
\[
 H_{XYZ}=\sum_{j=1}^N\big(J_x\sigma^x_j\sigma^x_{j+1}+J_y\sigma^y_j\sigma^y_{j+1}+J_z\sigma^z_j\sigma^z_{j+1}\big),
 \qquad J_x:J_y:J_z=1:\Gamma:\Delta,
\]
is, up to the additive constant $N(a'+c')/(2a(0))$ and the factor
$2a(0)/(b'+d')$, the logarithmic derivative of $t(u)$ at $u=0$; hence
$[H_{XYZ},t(u)]=0$.
\item Hierarchy.  Show: $p\to0$ gives $\Gamma=1$, $J_x=J_y$, the XXZ chain
with $\Delta=\cos\eta$ (after $\eta\to\ii\eta$, the $\Delta=\cosh\eta$
form of Part V); the further limit $\eta\to0$ with $u/\eta$ fixed gives the
XXX chain and the Yang $R$-matrix; and the limits $k\to1$ of the modulus
(via $\tau\mapsto-1/\tau$) give the hyperbolic (massive XXZ) forms.  State
the resulting inclusions rational $\subset$ trigonometric $\subset$
elliptic of solutions of the Yang--Baxter equation.
\end{enumerate}
\end{exercise}

\section{Fusion}

\begin{exercise}[The fusion procedure: degeneration of $R$ at $u=-\eta$ and the fused $R$-matrix on $V_2\otimes V_1$]
Let $R(u)$ be either the six-vertex $R$-matrix of Part IV or the elliptic
eight-vertex $R$-matrix, both with $R(0)\propto P$.  Write
$P^{\pm}=\frac12(1\pm P)$ for the projections onto the symmetric and
antisymmetric parts of $V\otimes V$.
\begin{enumerate}[label=(\alph*)]
\item Prove $R(-\eta)=2b(-\eta)P^-$ in both cases (six-vertex:
$a(-\eta)=0$, $c(-\eta)=-b(-\eta)$; eight-vertex: $a(-\eta)=d(-\eta)=0$,
$c(-\eta)=-b(-\eta)$).
\item From the Yang--Baxter equation at $u-v=-\eta$ deduce
$P^-_{12}R_{13}(u)R_{23}(u+\eta)=R_{23}(u+\eta)R_{13}(u)P^-_{12}$ and hence
that $R_{13}(u)R_{23}(u+\eta)$ maps $(\operatorname{im}P^+_{12})\otimes V$
into itself and $R_{23}(u+\eta)R_{13}(u)$ maps $(\operatorname{im}P^-_{12})\otimes V$
into itself.
\item Define the fused $R$-matrix
$R^{(2,1)}(u)=R_{13}(u)R_{23}(u+\eta)\big|_{(\operatorname{im}P^+_{12})\otimes V}\in\End(\Sym^2V\otimes V)$
and prove the mixed Yang--Baxter equation
\[
 R^{(2,1)}_{A3}(u-v)R^{(2,1)}_{A4}(u)R_{34}(v)=R_{34}(v)R^{(2,1)}_{A4}(u)R^{(2,1)}_{A3}(u-v)
\]
on $\Sym^2V\otimes V\otimes V$ ($A$ the fused pair of factors $1,2$), by
expanding each fused matrix, commuting factors on disjoint spaces, and
applying the Yang--Baxter equation twice.
\item Quantum determinant.  Prove that $R_{23}(u+\eta)R_{13}(u)$ acts on
$(\operatorname{im}P^-_{12})\otimes V$ as the scalar
$\delta(u)=b(u)a(u+\eta)-d(u)d(u+\eta)=a(u)b(u+\eta)-c(u)c(u+\eta)$
(equal expressions by the identities of the eight-vertex section; for the
six-vertex case $\delta(u)=\sinh u\,\sinh(u+2\eta)$).
\end{enumerate}
\end{exercise}

\begin{exercise}[The fusion hierarchy of transfer matrices and the first functional relation]
Let $\mathcal T(u)=L_N(u)\cdots L_1(u)$ with $L_j(u)=R_{0j}(u)$ and
$t(u)=\Tr_0\mathcal T(u)$.  Define the fused monodromy
$\mathcal T^{(2)}(u)=\mathcal T_{0}(u)\mathcal T_{0'}(u+\eta)\big|_{\operatorname{im}P^+_{00'}\otimes\mathcal H_N}$
and $t^{(2)}(u)=\Tr_{\Sym^2V}\mathcal T^{(2)}(u)$.
\begin{enumerate}[label=(\alph*)]
\item Prove that $\mathcal T_0(u)\mathcal T_{0'}(u+\eta)=\prod_{j=N}^{1}\big(L_{j,0}(u)L_{j,0'}(u+\eta)\big)$
(factors with distinct $j$ commute), that
$P^-_{00'}\mathcal T_0(u)\mathcal T_{0'}(u+\eta)=\mathcal T_{0'}(u+\eta)\mathcal T_0(u)P^-_{00'}$,
and that $\mathcal T_{0'}(u+\eta)\mathcal T_0(u)$ acts on
$\operatorname{im}P^-_{00'}\otimes\mathcal H_N$ as the scalar $\delta(u)^N$.
\item Prove the functional relation
\[
 t(u)\,t(u+\eta)=t^{(2)}(u)+\delta(u)^N,
\]
by decomposing the trace over $V_0\otimes V_{0'}$ according to
$\operatorname{im}P^+\oplus\operatorname{im}P^-$ and using that the
operator $\mathcal T_0(u)\mathcal T_{0'}(u+\eta)$ is block triangular for
this decomposition.
\item Deduce $[t^{(2)}(u),t(v)]=0$ and
$[t^{(2)}(u),t^{(2)}(v)]=0$. If a common eigenspace of the
$t(u)$ has eigenvalue function $\Lambda(u)$, compute the
eigenvalue of $t^{(2)}(u)$ as
$\Lambda(u)\Lambda(u+\eta)-\delta(u)^N$.
\end{enumerate}
\end{exercise}

\section{Quantum dimensions}

\begin{exercise}[Restricted fusion and quantum dimensions of RSOS heights]
Put $\gamma=\pi/(k+2)$, $q=e^{\ii\gamma}$, and
$d_r=\sin((r+1)\gamma)/\sin\gamma$ for $0\le r\le k$.
These are the quantum dimensions of the height labels $r+1$;
they are the specialization of the weighted traces
$\operatorname{tr}_q(\id_{V_r})$ of the section
$U_q(\mathfrak{sl}_2)$.
Define $U_0(x)=1$, $U_1(x)=2x$,
$U_{r+1}(x)=2xU_r(x)-U_{r-1}(x)$ and
$N_r=U_r(A/2)$.
\begin{enumerate}[label=(\alph*)]
\item Show that $\psi(a)=\sin(a\gamma)$ satisfies
$A\psi=2\cos\gamma\,\psi$, that $d_0=1$, and
$d_1d_r=d_{r-1}+d_{r+1}$ with $d_{-1}=d_{k+1}=0$.
\item Define
$S_{ar}=\sqrt{2/(k+2)}\sin(a(r+1)\gamma)$,
$1\le a\le k+1$, $0\le r\le k$.
Sum geometric series to prove
$\sum_aS_{ar}S_{as}=\delta_{rs}$ and the inversion
$\sum_rS_{ar}S_{br}=\delta_{ab}$.
Deduce the kernel resolution
\[
 (A^N)_{ab}=\sum_{r=0}^k S_{ar}S_{br}
                   \bigl(2\cos((r+1)\gamma)\bigr)^N.
\]
\item Use the Chebyshev recurrence and $U_{k+1}(A/2)=0$ to prove
\[
 N_rN_s=\sum_{\substack{t=|r-s|\\t\equiv r+s\ (2)}}^{
                  \min(r+s,2k-r-s)}N_t.
\]
Show by induction that all $N_r$ have nonnegative integer
entries and that evaluation on $\psi$ gives
$d_rd_s=\sum_tN_{rs}^t d_t$ with the displayed fusion coefficients.
\end{enumerate}
\end{exercise}

\begin{exercise}[Critical RSOS face weights and the Temperley--Lieb identity]
For an admissible height path $(a_0,\dots,a_N)$ define
\[
 (e_i f)(a)=\mathbf1_{a_{i-1}=a_{i+1}}
 \sum_{b:\,|b-a_{i-1}|=1}
 \frac{\sqrt{\psi(a_i)\psi(b)}}{\psi(a_{i-1})}
 f(a_0,\dots,a_{i-1},b,a_{i+1},\dots,a_N),
\]
with $\psi$ from the preceding exercise and $1\le i<N$.
\begin{enumerate}[label=(\alph*)]
\item Prove $e_i^*=e_i$, $e_i^2=d_1e_i$,
$e_ie_{i+1}e_i=e_i$, and distant commutativity.
Identify $e_i/d_1$ as a rank-one orthogonal projection on
each returning-path fibre.
\item Put $X_i(u)=\sin(\gamma-u)\id+\sin u\,e_i$.
Prove positivity of its admissible kernel for $0<u<\gamma$
and derive the face Yang--Baxter equation from the
Baxterization exercise of the section Hecke / Temperley--Lieb
algebras by analytic continuation.
\end{enumerate}
\end{exercise}

\section{Elliptic quantum groups}

\begin{exercise}[A dynamical elliptic vertex operator]
Write $\theta=\vartheta_1(\,\cdot\,|\tau)$ and
$s(a,z)=\theta(a-z)\theta'(0)/(\theta(a)\theta(z))$.
For $\gamma\ne0$ define a meromorphic arrow kernel
$\mathcal R(u,\lambda)$ equal to one on $++$ and $--$,
with mixed-arrow block
\[
 \frac1{s(\gamma,u)}
 \begin{pmatrix}s(\gamma,\lambda)&s(\lambda,u)\\
                 s(-\lambda,u)&s(\gamma,-\lambda)\end{pmatrix}.
\]
Set $\epsilon=\gamma/2$, and interpret
$f(\lambda+\epsilon Z_j)$ through the two spectral projections
of the physical arrow involution $Z_j$.
\begin{enumerate}[label=(\alph*)]
\item Use the theta addition identity to verify
\[
\begin{split}
 &\mathcal R_{12}(u-v,\lambda+\epsilon Z_3)
 \mathcal R_{13}(u,\lambda-\epsilon Z_2)
 \mathcal R_{23}(v,\lambda+\epsilon Z_1)\\
 &\quad=\mathcal R_{23}(v,\lambda-\epsilon Z_1)
 \mathcal R_{13}(u,\lambda+\epsilon Z_2)
 \mathcal R_{12}(u-v,\lambda-\epsilon Z_3).
\end{split}
\]
Work away from poles and extend meromorphically.
\item Define a representation of this elliptic exchange
relation to be a weight space $W$ and an operator $L(u,\lambda)$
satisfying the same equation with the last factor replaced by
$W$ and $\mathcal R_{13},\mathcal R_{23}$ replaced by $L_{13},L_{23}$.
Verify that $W=V$, $L(u,\lambda)=\mathcal R(u-a,\lambda)$
is such a representation for each $a$.
Compare conventions with Felder \cite{felder-elliptic}.
\end{enumerate}
\end{exercise}

\begin{exercise}[Prediction: residual entropy of a constrained RSOS strip]
Use the uniform RSOS chain of the section Theta and elliptic
functions, with $N$ bonds, fixed initial height $a$, and freely
summed terminal height. Let
$Z_N(a)=\sum_b(A^N)_{ab}$, and define the residual entropy
per bond as $s_0=k_B\lim_{N\to\infty}N^{-1}\log Z_N(a)$.
\begin{enumerate}[label=(\alph*)]
\item Since $\psi$ is strictly positive, bound $Z_N(a)$
above and below by positive constants times
$(2\cos\gamma)^N\psi(a)$, using
$\sum_b(A^N)_{ab}\psi(b)=(2\cos\gamma)^N\psi(a)$.
Derive
\[
 s_0=k_B\log\bigl(2\cos(\pi/(k+2))\bigr).
\]
Explain why summing the endpoint avoids the parity zeros
of a fixed-endpoint partition function.
\item Obtain $s_0=0$ for $k=1$, $s_0=\frac12k_B\log2$
for $k=2$, and $s_0=k_B\log((1+\sqrt5)/2)$ for $k=3$.
Derive the physical free energy per bond $-Ts_0$ for
the exactly degenerate allowed configurations.
\item Use the spectral kernel to calculate finite-$N$
corrections. Specify that an entropy measurement tests
this counting law only when forbidden neighbouring heights
cost energy much larger than $k_BT$ and energy splittings
among allowed configurations are much smaller than $k_BT$.
\end{enumerate}
\end{exercise}
